\documentclass[11pt]{article}

\usepackage[T1]{fontenc}
\usepackage{lmodern}
\usepackage[margin=0.78in]{geometry}
\usepackage{amsmath,amssymb,mathtools}
\usepackage{booktabs,enumitem}
\usepackage{xcolor}
\usepackage{microtype}
\usepackage[colorlinks=true,linkcolor=blue!45!black]{hyperref}
\usepackage{fancyhdr}

\setlength{\parindent}{1.05em}
\setlength{\parskip}{2pt}
\setlength{\abovedisplayskip}{7pt}
\setlength{\belowdisplayskip}{7pt}
\setlength{\abovedisplayshortskip}{4pt}
\setlength{\belowdisplayshortskip}{4pt}
\raggedbottom

\pagestyle{fancy}
\fancyhf{}
\fancyhead[L]{\footnotesize Accepted-refit Fubini--Study whitening}
\fancyhead[R]{\footnotesize RA-ADAPT}
\fancyfoot[C]{\thepage}
\renewcommand{\headrulewidth}{0.25pt}

\newcommand{\ket}[1]{\lvert #1\rangle}
\newcommand{\bra}[1]{\langle #1\rvert}
\newcommand{\braket}[2]{\langle #1\mid #2\rangle}
\newcommand{\R}{\mathbb R}
\newcommand{\conceptbreak}{%
  \par\medskip\noindent
  {\color{black!30}\rule{\textwidth}{0.35pt}}%
  \par\medskip}
\newcommand{\commentary}{\paragraph{Commentary.}}
\newcommand{\physical}{\paragraph{Physical meaning.}}
\newcommand{\paperone}{\paragraph{Paper-I consequence.}}

\title{Accepted-Refit Fubini--Study Whitening in RA-ADAPT\\
\large A Paper-I notation companion for the fixed Powell chart}
\author{}
\date{29 July 2026}

\begin{document}
\maketitle
\vspace{-1.25em}

\begin{abstract}
This companion explains the Fubini--Study whitening used only for the inner
optimization after RA-ADAPT admits a candidate.  The construction begins with
the accepted ansatz's logical amplitudes
\(\boldsymbol\theta\), maps logical displacements into physical horizontal
tangent motion, and then defines a supported matrix \(W\) whose columns are
orthonormal in the accepted Fubini--Study Gram metric.  Powell optimizes the
resulting coordinates \(\mathbf y\) through the fixed chart
\(\boldsymbol\theta(\mathbf y)
=\boldsymbol\theta^{(0)}+W\mathbf y\).
The expanded runtime coordinates \(\mathbf x\) and the block-averaging map
\(\mathsf B\) are treated afterward as one implementation realization of that
abstract logical-coordinate construction.  They are not additional physical
generator amplitudes.  Retained support appears only to remove Gram-null
directions and make the inverse square root well defined.  No candidate
response model or trust-region solve is derived here.
\end{abstract}

\noindent\fcolorbox{black!30}{blue!3}{%
\parbox{0.97\textwidth}{\small
\textbf{The entire subject of this companion is the accepted refit.}
A candidate has already been selected and inserted.  The only question is:
how are the enlarged ansatz amplitudes parameterized for the ensuing Powell
optimization?  The answer is a fixed, supported, locally Fubini--Study
orthonormal chart.}}

\section*{1. Admission enlarges the chart without changing the state}

At outer iteration \(k\), let the optimized ansatz contain \(n_k\) active
logical amplitudes
\(\boldsymbol\theta_k^\star\).
After a singleton candidate is admitted, the enlarged ansatz contains
\begin{equation}
d:=n_k+1
\label{eq:logical-dimension}
\end{equation}
logical amplitudes.  Its initial post-admission coordinate vector is
\(\boldsymbol\theta^{(0)}\in\R^d\): the previously optimized amplitudes appear
in their post-admission ordered positions, and the newly admitted amplitude
is zero.  Therefore
\begin{equation}
\ket{\psi_{k+1}(\boldsymbol\theta^{(0)})}
=
\ket{\psi_k(\boldsymbol\theta_k^\star)}.
\label{eq:zero-growth}
\end{equation}

Schematically, the accepted circuit consists of ordered logical generator
blocks,
\begin{equation}
U_{k+1}(\boldsymbol\theta)
=
U_d(\theta_d)\cdots U_2(\theta_2)U_1(\theta_1),
\qquad
U_\ell(\theta_\ell)=e^{-i\theta_\ell A_\ell}.
\label{eq:logical-block-circuit}
\end{equation}
The scalar \(\theta_\ell\) is the single variational amplitude controlling
logical generator block \(\ell\).  A block can later compile into several
executable factors, but that compilation does not create several independent
logical amplitudes for the generator.

Three logical-coordinate symbols must remain distinct:
\begin{center}
\renewcommand{\arraystretch}{1.18}
\begin{tabular}{@{}cl@{}}
\toprule
Symbol & Meaning \\
\midrule
\(\boldsymbol\theta^{(0)}\) &
inherited post-admission starting point \\
\(\boldsymbol\theta\) &
generic logical amplitudes during the accepted refit \\
\(\boldsymbol\theta_{k+1}^\star\) &
optimized logical endpoint returned by the refit \\
\bottomrule
\end{tabular}
\end{center}

\commentary Equation~\eqref{eq:zero-growth} says that admission changes the
coordinate description before it changes the prepared state.  The enlarged
circuit contains one more rotation, but that rotation is initially the
identity.  Nevertheless, changing its amplitude away from zero can produce a
nonzero physical tangent direction.  The accepted refit is precisely the
optimization that decides how far to move along the new direction and how
the previously active amplitudes should readjust.

\physical The state is the same at the instant of admission, but the local
space of available state motions can be larger.  Whitening is constructed at
this common physical anchor.  It calibrates the coordinate axes of the
enlarged chart before Powell begins moving away from that anchor.

\noindent\fcolorbox{black!30}{yellow!8}{%
\parbox{0.97\textwidth}{\small
\textbf{Concrete reading.}
If the accepted ansatz contains five generator blocks, then
\(\boldsymbol\theta\in\mathbb R^5\).  The entry \(\theta_3\) is the amplitude
of the third block in the ordered circuit.  It is not an eigenmode, a Pauli
word label, or a Powell direction.  It is the ordinary logical circuit
amplitude that the final map
\(\boldsymbol\theta(\mathbf y)=\boldsymbol\theta^{(0)}+W\mathbf y\)
must produce.}}

\section*{2. What is actually being whitened?}

Whitening does not act directly on a Hamiltonian, an energy Hessian, or a
candidate score.  It acts on the coordinate representation of first-order
physical state motion generated by the admitted ansatz.

At \(\boldsymbol\theta^{(0)}\), define
\begin{equation}
\ket{\partial_\ell\psi}_{0}
:=
\left.
\frac{\partial}{\partial\theta_\ell}
\ket{\psi_{k+1}(\boldsymbol\theta)}
\right|_{\boldsymbol\theta=\boldsymbol\theta^{(0)}},
\qquad
\ell=1,\ldots,d.
\label{eq:logical-derivatives}
\end{equation}
If \(\mathbf e_\ell\) is the \(\ell\)th ordinary basis vector in
\(\mathbb R^d\), this derivative means
\begin{equation}
\ket{\partial_\ell\psi}_{0}
=
\lim_{\epsilon\to0}
\frac{
\ket{\psi_{k+1}(\boldsymbol\theta^{(0)}
+\epsilon\mathbf e_\ell)}
-
\ket{\psi_{k+1}(\boldsymbol\theta^{(0)})}
}{\epsilon}.
\label{eq:logical-derivative-limit}
\end{equation}
In words: vary only logical amplitude \(\theta_\ell\), hold every other
accepted amplitude fixed, and record the state change per unit amplitude.

Paper I removes the unphysical global-phase component with
\begin{equation}
\Pi
:=
\mathbb I
-
\ket{\psi_{k+1}(\boldsymbol\theta^{(0)})}
\bra{\psi_{k+1}(\boldsymbol\theta^{(0)})}.
\label{eq:projector}
\end{equation}
A small simultaneous logical displacement obeys the first-order expansion
\begin{equation}
\ket{\psi_{k+1}(\boldsymbol\theta^{(0)}
+\delta\boldsymbol\theta)}
=
\ket{\psi_{k+1}(\boldsymbol\theta^{(0)})}
+
\sum_{\ell=1}^{d}
\delta\theta_\ell\ket{\partial_\ell\psi}_{0}
+
O(\|\delta\boldsymbol\theta\|_2^2).
\label{eq:state-first-order-expansion}
\end{equation}
This is why several derivative terms appear: a general displacement can
change several accepted amplitudes at once, so its first-order state motion is
the corresponding linear combination of their individual motions.

A logical parameter displacement
\(\delta\boldsymbol\theta\in\R^d\) therefore produces the horizontal physical
tangent
\begin{equation}
\ket{\delta\psi_\perp}
:=
\Pi
\sum_{\ell=1}^{d}
\delta\theta_\ell\ket{\partial_\ell\psi}_{0}.
\label{eq:physical-tangent-map}
\end{equation}

Read Eq.~\eqref{eq:physical-tangent-map} one symbol at a time:
\begin{center}
\renewcommand{\arraystretch}{1.18}
\begin{tabular}{@{}p{0.20\textwidth}p{0.73\textwidth}@{}}
\toprule
Object & Literal meaning \\
\midrule
\(\delta\theta_\ell\) &
small numerical change of accepted logical amplitude \(\ell\) \\
\(\ket{\partial_\ell\psi}_{0}\) &
state change per unit change of that one amplitude at the origin \\
\(\sum_\ell\delta\theta_\ell\ket{\partial_\ell\psi}_{0}\) &
total first-order state change produced by changing all selected amplitudes \\
\(\Pi\) &
removes the component that changes only the global phase of the state \\
\(\ket{\delta\psi_\perp}\) &
remaining physically distinguishable first-order motion of the quantum ray \\
\bottomrule
\end{tabular}
\end{center}

This is the central map:
\begin{equation}
\boxed{
\delta\boldsymbol\theta
\quad\longmapsto\quad
\ket{\delta\psi_\perp}.
}
\label{eq:logical-to-physical}
\end{equation}
The vector \(\delta\boldsymbol\theta\) lives in a parameter-coordinate space.
The ket \(\ket{\delta\psi_\perp}\) lives in the physical tangent space of the
quantum ray at the accepted state.  These are not the same mathematical
object.

The accepted logical-coordinate Gram matrix is
\begin{equation}
G_{\ell q}^{(0)}
:=
\operatorname{Re}
\left[
{}_{0}\!\bra{\partial_\ell\psi}
\Pi
\ket{\partial_q\psi}_{0}
\right],
\qquad
\mathbf G^{(0)}\in\R^{d\times d}.
\label{eq:logical-gram}
\end{equation}
It is the pullback of the physical Fubini--Study inner product to the logical
parameter coordinates.  Consequently,
\begin{equation}
\left\|\ket{\delta\psi_\perp}\right\|_{\R}^{2}
=
\delta\boldsymbol\theta^\top
\mathbf G^{(0)}
\delta\boldsymbol\theta.
\label{eq:fs-length}
\end{equation}

\commentary The entries of \(\mathbf G^{(0)}\) answer two related questions.
The diagonal entry \(G_{\ell\ell}^{(0)}\) gives the squared first-order
Fubini--Study length produced by a unit change of logical amplitude
\(\theta_\ell\).  The off-diagonal entry \(G_{\ell q}^{(0)}\) records the
physical overlap between the motions produced by changing
\(\theta_\ell\) and \(\theta_q\).  Thus the Gram matrix measures both
coordinate scale and collective redundancy.

\physical An optimizer sees numerical coordinates.  The quantum state sees
physical tangent motion.  If one raw coordinate moves the state much farther
than another, equal numerical steps do not have equal physical meaning.  If
two coordinates produce nearly the same tangent, searching them independently
can waste optimizer effort.  Whitening uses \(\mathbf G^{(0)}\) to give the
optimizer axes a common local physical scale.

\subsection*{The edge cases encoded by the projector and Gram matrix}

\paragraph{Pure-phase direction.}
It is possible for a parameter change to alter only the global phase:
\[
\sum_\ell\delta\theta_\ell\ket{\partial_\ell\psi}_{0}
=
i c\ket{\psi_{k+1}(\boldsymbol\theta^{(0)})}.
\]
The projector sends this motion to zero.  Such a change does not move the
physical quantum ray and should not receive a Powell coordinate.

\paragraph{Redundant logical directions.}
Two different accepted generator blocks can produce the same horizontal
tangent at the present state.  Then changing them in opposite combinations
can give
\(\ket{\delta\psi_\perp}=\mathbf 0\), even though
\(\delta\boldsymbol\theta\ne\mathbf0\).  This is a physical tangent
redundancy of the accepted ansatz and appears as a null vector of
\(\mathbf G^{(0)}\).

\paragraph{Different physical scales.}
Two logical directions can be independent but have very different tangent
lengths.  If
\(G_{11}^{(0)}=100\) and \(G_{22}^{(0)}=1\), the same raw amplitude step moves
the state ten times farther along direction \(1\) than along direction \(2\).
Whitening rescales the directions so Powell does not inherit that arbitrary
coordinate imbalance.

\conceptbreak

\section*{3. Retained support: the one preliminary step}

The Gram matrix is positive semidefinite, not necessarily positive definite.
A vector
\(\mathbf a\in\ker\mathbf G^{(0)}\) satisfies
\begin{equation}
\mathbf a^\top\mathbf G^{(0)}\mathbf a=0,
\end{equation}
so the logical combination \(\delta\boldsymbol\theta=\mathbf a\) produces no
resolved first-order physical motion at the post-admission origin.

In the one-coordinate-per-generator description, diagonalize the
\emph{logical} Gram matrix directly.  Let \(r\) be the number of eigenmodes
retained by Paper I's relative support rule, let
\(V\in\R^{d\times r}\) collect their eigenvectors, and let
\(\Lambda\in\R^{r\times r}\) contain their positive eigenvalues:
\begin{equation}
\mathbf G^{(0)}V=V\Lambda,
\qquad
V^\top V=I_r,
\qquad
\Lambda\succ0.
\label{eq:logical-retained-eigensystem}
\end{equation}
The columns of \(V\) are collective combinations of the original logical
generator amplitudes.  They span the retained logical-coordinate tangent
support.

\commentary Support is not the main conceptual difficulty here.  It performs
one necessary piece of housekeeping: a zero or numerically unresolved Gram
eigenvalue cannot be inverse-square-rooted.  Removing those modes prevents
Powell from receiving an optimizer coordinate that produces no certified
first-order physical motion.

\physical The support reduction does not delete an accepted generator from
the circuit.  It removes a locally unresolved combination from the optimizer
chart constructed at this particular anchor.  When \(r=d\), all logical
directions are retained.  When \(r<d\), the fixed Powell chart is a supported
\(r\)-dimensional subchart of the \(d\)-dimensional logical parameter space.

\paragraph{Three cases.}
\begin{enumerate}[leftmargin=1.7em,itemsep=2pt,topsep=3pt]
\item If every accepted logical tangent is independent and resolved, then
\(r=d\); support changes no dimension.
\item If two accepted logical directions generate exactly the same physical
tangent, one difference combination is Gram-null and \(r<d\).
\item If a direction is not exactly null but has an eigenvalue below the
declared relative tolerance, it is treated as numerically unresolved.  This
guards the inverse square root against amplifying an almost-null direction
into an enormous optimizer step.
\end{enumerate}

\section*{4. Whitening with one coordinate per generator}

The inverse square root rescales each retained eigenvector by the inverse
square root of its physical length:
\begin{equation}
\boxed{
W:=V\Lambda^{-1/2}
\in\R^{d\times r}.
}
\label{eq:single-coordinate-W}
\end{equation}
Equivalently, write
\begin{equation}
W
=
\begin{pmatrix}
\mathbf w_1&\cdots&\mathbf w_r
\end{pmatrix}
\in\R^{d\times r},
\label{eq:W-columns}
\end{equation}
where each column is
\(\mathbf w_a=\mathbf v_a/\sqrt{\lambda_a}\).  These columns span the same
retained logical-coordinate support as \(V\), but are normalized so that
\begin{equation}
\boxed{
W^\top\mathbf G^{(0)}W=I_r.
}
\label{eq:logical-whitening}
\end{equation}
Indeed,
\begin{align}
W^\top\mathbf G^{(0)}W
&=
\Lambda^{-1/2}V^\top
\mathbf G^{(0)}
V\Lambda^{-1/2}\nonumber\\
&=
\Lambda^{-1/2}\Lambda\Lambda^{-1/2}
=I_r.
\label{eq:single-coordinate-whitening-proof}
\end{align}

The whitened optimizer coordinates are
\(\mathbf y\in\R^r\), and the accepted ansatz is presented to Powell through
\begin{equation}
\boxed{
\boldsymbol\theta(\mathbf y)
=
\boldsymbol\theta^{(0)}
+W\mathbf y.
}
\label{eq:powell-chart}
\end{equation}
Thus a whitened displacement
\(\delta\mathbf y\) produces the logical displacement
\begin{equation}
\delta\boldsymbol\theta=W\,\delta\mathbf y.
\label{eq:y-to-theta}
\end{equation}
Substituting Eq.~\eqref{eq:y-to-theta} into the Fubini--Study length gives
\begin{align}
\left\|\ket{\delta\psi_\perp}\right\|_{\R}^{2}
&=
\delta\boldsymbol\theta^\top
\mathbf G^{(0)}
\delta\boldsymbol\theta\nonumber\\
&=
\delta\mathbf y^\top
W^\top\mathbf G^{(0)}W
\delta\mathbf y\nonumber\\
&=
\delta\mathbf y^\top\delta\mathbf y.
\label{eq:whitened-fs-length}
\end{align}

Equation~\eqref{eq:whitened-fs-length} is the entire point of whitening.
Euclidean length in the optimizer coordinates equals first-order
Fubini--Study length at the accepted-refit origin.

\subsection*{The physical tangent frame defined by \(W\)}

For each optimizer coordinate \(y_a\), define the corresponding physical
tangent
\begin{equation}
\ket{e_a}
:=
\Pi
\sum_{\ell=1}^{d}
W_{\ell a}\ket{\partial_\ell\psi}_{0},
\qquad
a=1,\ldots,r.
\label{eq:whitened-tangents}
\end{equation}
Then Eq.~\eqref{eq:logical-whitening} is equivalent to
\begin{equation}
\operatorname{Re}\braket{e_a}{e_b}
=
\delta_{ab}.
\label{eq:orthonormal-tangents}
\end{equation}
Therefore:
\begin{enumerate}[leftmargin=1.6em,itemsep=2pt,topsep=3pt]
\item every \(y_a\) produces a retained physical tangent;
\item one unit of \(y_a\) produces one unit of first-order Fubini--Study
length;
\item distinct \(y_a\) and \(y_b\) produce Fubini--Study-orthogonal tangents
at the origin.
\end{enumerate}

\commentary The columns of \(W\) are not additional generators.  Each column
is a collective logical-amplitude direction.  It tells Powell which
combination of accepted circuit amplitudes to change when Powell moves along
one numerical axis.

For example, if
\[
\mathbf w_1
=
\begin{pmatrix}
1/2\\-1/4\\0
\end{pmatrix},
\]
then increasing \(y_1\) by a small amount \(\epsilon\) changes the accepted
logical amplitudes by
\[
\delta\theta_1=\epsilon/2,\qquad
\delta\theta_2=-\epsilon/4,\qquad
\delta\theta_3=0.
\]
The column is a recipe for a collective circuit move.  The normalization
\(\mathbf w_1^\top\mathbf G^{(0)}\mathbf w_1=1\) says that a unit numerical
move along that recipe has unit first-order Fubini--Study length.

\physical Whitening makes the \emph{kinematic} tangent frame locally
orthonormal.  It does not say that equal physical motions have equal energy
effects.  The Hamiltonian can still make one whitened direction steep, another
flat, and two directions energetically coupled.  Whitening standardizes how
far the state moves, not how the energy varies along that motion.

\section*{5. What whitening does and does not guarantee}

\subsection*{It is local}

The identity
\(W^\top\mathbf G^{(0)}W=I_r\) holds at
\(\boldsymbol\theta^{(0)}\), where \(\mathbf G^{(0)}\) was evaluated.  At a
finite Powell displacement, define
\begin{equation}
\mathbf G_y(\mathbf y)
:=
W^\top
\mathbf G[\boldsymbol\theta(\mathbf y)]
W.
\label{eq:evolving-metric}
\end{equation}
Then
\begin{equation}
\mathbf G_y(\mathbf 0)=I_r,
\end{equation}
but generally
\(\mathbf G_y(\mathbf y)\ne I_r\) for
\(\mathbf y\ne\mathbf 0\).

\commentary The chart is whitened once at the post-admission origin and then
held fixed.  It is a local preconditioner, not a continuous change of
coordinates recomputed at every Powell evaluation.

\subsection*{It does not make the energy landscape spherical}

Define the accepted-ansatz energy in logical coordinates by
\begin{equation}
E_{k+1}(\boldsymbol\theta)
:=
\bra{\psi_{k+1}(\boldsymbol\theta)}
H
\ket{\psi_{k+1}(\boldsymbol\theta)}.
\label{eq:logical-energy}
\end{equation}
Whitening composes this same scalar function with the affine map
\(\boldsymbol\theta(\mathbf y)\).  The resulting function can remain
anisotropic, curved, nonquadratic, and multimodal.  Only the local physical
lengths and angles of the coordinate tangents are normalized.

\subsection*{It is a reparameterization only on the retained chart}

If \(r=d\), \(W\) is square and the chart is a local reparameterization of all
logical directions.  If \(r<d\), Powell searches only the affine supported
subchart
\begin{equation}
\boldsymbol\theta^{(0)}
+\operatorname{range}(W).
\label{eq:supported-affine-chart}
\end{equation}
This explicitly excludes Gram-null or numerically unresolved combinations at
the construction point.

\paperone The main methodology should state this distinction plainly.  It is
too strong to describe a rank-deficient \(W\) as merely an invertible
coordinate relabeling of the entire accepted parameter space.  It is instead
a locally supported optimizer chart.

\conceptbreak

\section*{6. What Powell actually optimizes}

Powell does not receive \(\boldsymbol\theta\) as its numerical variable.  It
receives
\begin{equation}
\mathbf y\in\R^r
\end{equation}
and evaluates the exact accepted-ansatz energy
\begin{equation}
E_{k+1}(\mathbf y)
:=
\bra{\psi_{k+1}(\boldsymbol\theta(\mathbf y))}
H
\ket{\psi_{k+1}(\boldsymbol\theta(\mathbf y))}.
\label{eq:powell-energy}
\end{equation}
The optimization begins at
\(\mathbf y=\mathbf 0\), because
\begin{equation}
\boldsymbol\theta(\mathbf 0)
=
\boldsymbol\theta^{(0)}.
\end{equation}
In compact notation,
\begin{equation}
\mathbf y^\star
\in
\operatorname*{arg\,min}_{\mathbf y\in\R^r}
E_{k+1}(\mathbf y),
\qquad
\boldsymbol\theta_{k+1}^\star
=
\boldsymbol\theta^{(0)}+W\mathbf y^\star.
\label{eq:powell-endpoint}
\end{equation}

\commentary Powell constructs and updates numerical search directions using
the coordinates it is given.  With raw logical amplitudes, those numerical
directions can have very different local physical lengths and can overlap
strongly in the state tangent plane.  With whitened coordinates, Powell
begins from resolved axes having equal local Fubini--Study length and zero
pairwise Fubini--Study overlap.

\physical The state family and energy evaluations remain quantum-mechanical;
the whitening calculation is a classical construction of the coordinate
chart in which those evaluations are requested.  The optimizer explores
\(\mathbf y\), the chart converts \(\mathbf y\) into logical circuit
amplitudes, and the circuit evaluates the corresponding energy.

\subsection*{Why the chart remains fixed for one invocation}

The matrices \(\mathbf G^{(0)}\) and \(W\) are constructed once at the
post-admission origin.  Every objective evaluation in the ensuing Powell
refit uses the same map
\begin{equation}
\mathbf y
\longmapsto
\boldsymbol\theta^{(0)}+W\mathbf y.
\label{eq:fixed-chart-map}
\end{equation}
This gives all Powell line searches a common numerical coordinate system.
After the next candidate admission, a new enlarged ansatz, new origin, and
new accepted Gram matrix define the next whitening map.

\section*{7. Extension: multiple stored coordinates per generator}

Everything through Section~6 defines the mathematical whitening method in
logical ansatz coordinates.  This section explains one reported
implementation realization.  Its symbols need not appear in the main
methodology if they obscure the physical construction.

\subsection*{One generator still has one logical amplitude}

Let
\begin{equation}
\boldsymbol\theta\in\R^d
\end{equation}
contain one logical amplitude for each accepted generator block.  The
implementation can store several executable Pauli-factor parameters for one
logical generator, but those stored values are averaged before execution.
They are not additional independent logical generator amplitudes.

\subsection*{What an \(x_a\) actually represents}

A logical block can be realized by several executable factors.  During the
reported metric construction, the implementation exposes a temporary
coordinate slot for each such factor.  The scalar \(x_a\) is the value stored
in temporary slot \(a\); the vector
\begin{equation}
\mathbf x=(x_1,\ldots,x_m)^\top\in\R^m
\label{eq:x-definition}
\end{equation}
collects all \(m\) slots.

These slots are bookkeeping coordinates for the expanded executable
representation.  They are not \(m\) independent accepted-generator
amplitudes.  Before a logical block is interpreted, the slots belonging to
that block are averaged into its one logical amplitude displacement.

For example, suppose logical block \(\ell\) has three temporary slots
\(x_4,x_5,x_6\).  Around the fixed construction origin:
\begin{itemize}[leftmargin=1.7em,itemsep=2pt,topsep=3pt]
\item the synchronized slot displacement
\((\delta x_4,\delta x_5,\delta x_6)=(t,t,t)\) produces the logical
displacement \(\delta\theta_\ell=t\);
\item the disagreement displacement
\((t,-t,0)\) produces \(\delta\theta_\ell=0\), because its mean is zero;
\item Powell never receives \(x_4,x_5,x_6\).  After the whitening matrix has
been constructed, Powell receives only \(\mathbf y\), and \(W\mathbf y\)
directly produces logical amplitude displacements.
\end{itemize}

\noindent\fcolorbox{black!30}{yellow!8}{%
\parbox{0.97\textwidth}{\small
\textbf{Why \(\mathbf x\) exists.}
The metric-construction routine can work in an expanded
generator-exponential representation even though the accepted variational
circuit has one logical amplitude per generator block.  The vector
\(\mathbf x\) names that temporary expanded representation.
It is used to construct \(W\); it is not the variable optimized by Powell and
is not itself a vector in the physical quantum-state tangent plane.}}

\subsection*{The averaging map \(\mathsf B\)}

For logical generator \(\ell\in\{1,\ldots,d\}\), let
\begin{equation}
I_\ell\subseteq\{1,\ldots,m\}
\label{eq:block-index-set}
\end{equation}
be the set of temporary slot indices belonging to that one generator, and
let
\begin{equation}
b_\ell:=|I_\ell|
\label{eq:block-size}
\end{equation}
be the number of its slots.  Define
\(\mathsf B\in\R^{d\times m}\) by
\begin{equation}
\mathsf B_{\ell a}
=
\begin{cases}
b_\ell^{-1}, & a\in I_\ell,\\
0, & a\notin I_\ell.
\end{cases}
\label{eq:B-entrywise}
\end{equation}
Then
\begin{equation}
(\mathsf B\,\delta\mathbf x)_\ell
=
\frac{1}{b_\ell}
\sum_{a\in I_\ell}\delta x_a,
\qquad
\delta\boldsymbol\theta
=
\mathsf B\,\delta\mathbf x.
\label{eq:B-map}
\end{equation}

The \(\ell\)th row of \(\mathsf B\) therefore answers one literal question:
``which temporary slots belong to logical generator \(\ell\), and with what
averaging weight?''  Multiplying by \(\mathsf B\) does not yet differentiate
the quantum state.  It only turns a temporary stored displacement into a
logical circuit-amplitude displacement.

The complete physical chain is
\begin{equation}
\boxed{
\delta\mathbf x
\xmapsto{\ \mathsf B\ }
\delta\boldsymbol\theta
\longmapsto
\ket{\delta\psi_\perp}.
}
\label{eq:base-to-physical}
\end{equation}

\commentary A base-coordinate displacement is not itself a physical tangent
vector.  It lives in an implementation parameter space.  First
\(\mathsf B\) converts it into one displacement per logical accepted
generator.  Only the ansatz derivative and projector in
Eq.~\eqref{eq:physical-tangent-map} convert that logical displacement into a
physical horizontal tangent.

\paragraph{Important edge cases.}
\begin{enumerate}[leftmargin=1.7em,itemsep=2pt,topsep=3pt]
\item If every logical generator has exactly one temporary slot, then
\(m=d\), every \(b_\ell=1\), and \(\mathsf B=I_d\).  The expanded chart and
logical chart coincide, so \(\mathbf x\) adds nothing.
\item If several slots belong to one generator, any within-block disagreement
whose entries sum to zero lies in \(\ker\mathsf B\).  This is an
implementation-representation redundancy.
\item Even when \(\mathsf B=I_d\), two distinct logical generators can still
produce redundant physical tangents.  That is a state-dependent physical
redundancy recorded by \(\mathbf G^{(0)}\), not a storage redundancy created
by \(\mathbf x\).
\end{enumerate}

\subsection*{The metric in the expanded base chart}

Substituting
\(\delta\boldsymbol\theta=\mathsf B\,\delta\mathbf x\) into
Eq.~\eqref{eq:fs-length} gives
\begin{align}
\left\|\ket{\delta\psi_\perp}\right\|_{\R}^{2}
&=
\delta\mathbf x^\top
\mathsf B^\top\mathbf G^{(0)}\mathsf B
\delta\mathbf x,\nonumber\\
\mathbf G_x^{(0)}
&:=
\mathsf B^\top\mathbf G^{(0)}\mathsf B
\in\R^{m\times m}.
\label{eq:base-gram}
\end{align}
Let \(V_x\in\R^{m\times r}\) contain the retained eigenvectors of this
expanded-coordinate metric, and let
\(\Lambda_x\in\R^{r\times r}\) contain their positive eigenvalues:
\begin{equation}
\mathbf G_x^{(0)}V_x=V_x\Lambda_x.
\label{eq:base-eigensystem}
\end{equation}
The implementation realizes the logical whitening matrix as
\begin{equation}
\boxed{
W
=
\mathsf B V_x\Lambda_x^{-1/2}
\in\R^{d\times r}.
}
\label{eq:implemented-W}
\end{equation}
Indeed,
\begin{align}
W^\top\mathbf G^{(0)}W
&=
\Lambda_x^{-1/2}V_x^\top
\mathsf B^\top\mathbf G^{(0)}\mathsf B
V_x\Lambda_x^{-1/2}\nonumber\\
&=
\Lambda_x^{-1/2}V_x^\top
\mathbf G_x^{(0)}
V_x\Lambda_x^{-1/2}\nonumber\\
&=I_r.
\label{eq:implemented-whitening-proof}
\end{align}

In the logical-coordinate special case,
\begin{equation}
\mathbf x=\boldsymbol\theta,\qquad
m=d,\qquad
\mathsf B=I_d,
\label{eq:logical-special-case}
\end{equation}
and Eq.~\eqref{eq:implemented-W} reduces to the familiar supported inverse
square root
\(W=V\Lambda^{-1/2}\) of the logical Gram matrix.

\paperone The abstract method needs only
\(\mathbf G^{(0)}\), \(W\), and \(\mathbf y\), together with
\(W^\top\mathbf G^{(0)}W=I_r\).
The expanded coordinates \(\mathbf x\), block sets \(I_\ell\), block sizes
\(b_\ell\), and averaging map \(\mathsf B\) describe how one implementation
constructs \(W\).  They are more naturally placed in a numerical appendix.

\conceptbreak

\section*{8. One worked example}

This example shows, in one calculation, the distinction among logical
amplitudes, expanded coordinates, retained support, and Powell coordinates.

\subsection*{Step 1: logical amplitudes and their physical metric}

Suppose the accepted ansatz has two logical amplitudes
\[
\boldsymbol\theta
=
\begin{pmatrix}\theta_1\\\theta_2\end{pmatrix},
\qquad d=2,
\]
and at \(\boldsymbol\theta^{(0)}\) its logical Gram matrix is
\begin{equation}
\mathbf G^{(0)}
=
\begin{bmatrix}
4&0\\
0&1
\end{bmatrix}.
\label{eq:example-logical-gram}
\end{equation}
The squared first-order physical length of a logical displacement is
\begin{equation}
\left\|\ket{\delta\psi_\perp}\right\|_{\R}^{2}
=
4\,\delta\theta_1^2+\delta\theta_2^2.
\label{eq:example-logical-length}
\end{equation}
A raw unit change of \(\theta_1\) moves the state twice as far in
Fubini--Study length as a raw unit change of \(\theta_2\).

\subsection*{Step 2: retain and whiten the two logical coordinates}

There is one coordinate per generator, so no expanded implementation chart is
needed.  The logical eigensystem is simply
\begin{equation}
V=I_2,
\qquad
\Lambda
=
\begin{bmatrix}
4&0\\
0&1
\end{bmatrix}.
\label{eq:example-logical-eigensystem}
\end{equation}
Therefore
\begin{equation}
W
=
V\Lambda^{-1/2}
=
\begin{bmatrix}
\tfrac12&0\\
0&1
\end{bmatrix},
\qquad
W^\top\mathbf G^{(0)}W=I_2.
\label{eq:example-logical-W}
\end{equation}
Powell's coordinates act through
\begin{equation}
\begin{pmatrix}
\theta_1(\mathbf y)\\
\theta_2(\mathbf y)
\end{pmatrix}
=
\begin{pmatrix}
\theta_1^{(0)}\\
\theta_2^{(0)}
\end{pmatrix}
+
\begin{bmatrix}
\tfrac12&0\\
0&1
\end{bmatrix}
\begin{pmatrix}y_1\\y_2\end{pmatrix}.
\label{eq:example-powell-chart}
\end{equation}
A unit \(y_1\) move changes \(\theta_1\) by \(1/2\), producing squared
Fubini--Study length
\[
4(1/2)^2=1.
\]
A unit \(y_2\) move changes \(\theta_2\) by \(1\), also producing squared
Fubini--Study length
\[
1(1)^2=1.
\]
The complete whitening method is already finished.  Each generator had one
logical coordinate, and the two Powell axes now have equal first-order
physical length.

\subsection*{Extension of the same example: multiple stored coordinates}

Only now suppose the first logical generator is represented by two stored
executable factors and the second by one:
\begin{equation}
m=3,\qquad
I_1=\{1,2\},\qquad
I_2=\{3\}.
\end{equation}
Then
\begin{equation}
\mathsf B
=
\begin{bmatrix}
\tfrac12&\tfrac12&0\\
0&0&1
\end{bmatrix},
\qquad
\begin{bmatrix}
\delta\theta_1\\
\delta\theta_2
\end{bmatrix}
=
\begin{bmatrix}
(\delta x_1+\delta x_2)/2\\
\delta x_3
\end{bmatrix}.
\label{eq:example-B}
\end{equation}
There are three stored base coordinates but still only two logical generator
amplitudes.

The expanded-coordinate Gram matrix is
\begin{equation}
\mathbf G_x^{(0)}
=
\mathsf B^\top\mathbf G^{(0)}\mathsf B
=
\begin{bmatrix}
1&1&0\\
1&1&0\\
0&0&1
\end{bmatrix}.
\label{eq:example-base-gram}
\end{equation}
Its eigenvalues are
\[
2,\qquad 1,\qquad 0.
\]
The zero eigenvector is proportional to
\((1,-1,0)^\top\).  It changes \(x_1\) and \(x_2\) oppositely, leaving their
mean unchanged:
\[
\mathsf B
\begin{pmatrix}1\\-1\\0\end{pmatrix}
=
\begin{pmatrix}0\\0\end{pmatrix}.
\]
The null direction is a redundancy of the expanded storage representation,
not an additional physical generator direction.

Choose the retained expanded-coordinate eigensystem
\begin{equation}
V_x
=
\begin{bmatrix}
1/\sqrt2&0\\
1/\sqrt2&0\\
0&1
\end{bmatrix},
\qquad
\Lambda_x
=
\begin{bmatrix}
2&0\\
0&1
\end{bmatrix}.
\label{eq:example-expanded-eigensystem}
\end{equation}
Then
\begin{align}
W
&=
\mathsf B V_x\Lambda_x^{-1/2}\nonumber\\
&=
\begin{bmatrix}
\tfrac12&0\\
0&1
\end{bmatrix}.
\label{eq:example-expanded-W}
\end{align}
This is exactly the logical-coordinate whitening matrix already obtained in
Eq.~\eqref{eq:example-logical-W}.  The multi-coordinate construction does not
define a different physical method; it recovers the same logical \(W\) from a
more redundant implementation representation.

\commentary The worked example must be read in this order.  First, whitening
rescales one coordinate per accepted generator.  That is the mathematical
method.  Second, the expanded-coordinate extension shows how duplicated
storage slots can introduce an additional null disagreement direction before
recovering the same logical whitening map.

\section*{9. How this companion should guide Paper I}

\subsection*{Main methodology: state the abstract logical construction}

The reader-facing methodology needs only the inherited logical origin
\(\boldsymbol\theta^{(0)}\), its accepted logical Gram matrix
\(\mathbf G^{(0)}\), the retained logical eigensystem \(V,\Lambda\), the
whitening map \(W\), and Powell's coordinates \(\mathbf y\).  A concise
Paper-I statement is:
\begin{quote}
Let \(\mathbf G^{(0)}\in
\mathbb R^{(n_k+1)\times(n_k+1)}\) be the Fubini--Study Gram matrix of the
accepted ansatz at \(\boldsymbol\theta^{(0)}\), and let \(r\) be the number of
eigenmodes retained by the relative support rule.  Let
\(\mathbf G^{(0)}V=V\Lambda\) be the retained logical-coordinate
eigensystem and define
\(W:=V\Lambda^{-1/2}\in
\mathbb R^{(n_k+1)\times r}\).  Then
\(W^\top\mathbf G^{(0)}W=I_r\), and the accepted refit uses the fixed
whitened chart
\[
\mathbf y\in\mathbb R^r,\qquad
\boldsymbol\theta(\mathbf y)
=\boldsymbol\theta^{(0)}+W\mathbf y.
\]
\end{quote}

This is the method.  It says what Powell optimizes and what physical
normalization the chart satisfies without introducing runtime storage
details.

\subsection*{Numerical appendix: document the reported realization}

The numerical appendix can then explain that the reported implementation
constructs \(W\) through an expanded base chart:
\begin{equation}
\delta\boldsymbol\theta=\mathsf B\,\delta\mathbf x,
\qquad
\mathbf G_x^{(0)}
=\mathsf B^\top\mathbf G^{(0)}\mathsf B,
\qquad
W=\mathsf B V_x\Lambda_x^{-1/2}.
\label{eq:appendix-realization}
\end{equation}
That appendix is also the appropriate place to define the block sets
\(I_\ell\), block sizes \(b_\ell\), and mean-displacement convention.

\paperone The main text should explain the accepted logical amplitudes, their
physical Fubini--Study tangent metric, the supported \(W\), and Powell's fixed
\(\mathbf y\)-chart.  The numerical appendix should explain how expanded
runtime values are averaged and used to construct \(W\).  This boundary
prevents implementation storage coordinates from being mistaken for
independent variational generator amplitudes.

\section*{10. Final notation map}

\begin{center}
\renewcommand{\arraystretch}{1.2}
\begin{tabular}{@{}p{0.17\textwidth}p{0.76\textwidth}@{}}
\toprule
Object & Meaning \\
\midrule
\(\boldsymbol\theta^{(0)}\) &
inherited logical amplitudes immediately after zero-amplitude admission \\
\(\boldsymbol\theta\) &
generic logical amplitudes of the accepted ansatz \\
\(\boldsymbol\theta_{k+1}^\star\) &
logical endpoint returned by the accepted refit \\
\(\ket{\delta\psi_\perp}\) &
physical horizontal tangent produced by a logical displacement \\
\(\mathbf G^{(0)}\) &
logical-coordinate pullback of the Fubini--Study metric at the refit origin \\
\(r\) &
number of retained physical tangent modes \\
\(V,\Lambda\) &
retained eigensystem of the one-coordinate-per-generator logical Gram matrix \\
\(W\) &
logical whitening map \(V\Lambda^{-1/2}\) from optimizer coordinates to
accepted-amplitude displacements \\
\(\mathbf y\) &
the \(r\)-component coordinate vector optimized by Powell \\
\(\mathbf x\) &
expanded implementation-coordinate representation, not a physical tangent \\
\(\mathsf B\) &
implementation map averaging expanded block displacements into logical ones \\
\(\mathbf G_x^{(0)}\) &
accepted Gram metric expressed in the expanded base chart \\
\(V_x,\Lambda_x\) &
retained expanded-coordinate eigensystem used to recover the logical \(W\) \\
\bottomrule
\end{tabular}
\end{center}

\section*{The chain to remember}

\begin{equation}
\boxed{
\begin{gathered}
\mathbf G^{(0)}V=V\Lambda,\qquad
W=V\Lambda^{-1/2},\\[0.3em]
\mathbf y
\xmapsto[\text{fixed Powell chart}]{W}
\delta\boldsymbol\theta
\xmapsto[\text{ansatz derivative}]{}
\ket{\delta\psi_\perp},\\[0.3em]
\boldsymbol\theta(\mathbf y)
=
\boldsymbol\theta^{(0)}+W\mathbf y,
\qquad
\boldsymbol\theta_{k+1}^\star
=
\boldsymbol\theta^{(0)}+W\mathbf y^\star,\\[0.3em]
\text{extension only:}\quad
\delta\mathbf x\xmapsto{\mathsf B}\delta\boldsymbol\theta,\qquad
W=\mathsf B V_x\Lambda_x^{-1/2}.
\end{gathered}}
\label{eq:final-chain}
\end{equation}

\noindent
The first three lines give the complete one-coordinate-per-generator method.
The last line is the later implementation extension for a generator represented
by several temporary executable-coordinate slots.

\end{document}
